\section{Definición y futuro del Agent}

\subsection{Las dos características centrales del Agent}

\begin{importantbox}{Agent = multiturno + herramientas}
La definición que da Yang Zhilin (杨植麟) del Agent es concisa y esencial:
\begin{enumerate}
    \item \textbf{Multiturno} (Multi-turn): puede sostener interacciones múltiples; es una forma de Test-Time Scaling que permite al modelo completar tareas complejas
    \item \textbf{Herramientas} (Tool Use): la manera de conectar el cerebro con el mundo exterior; el motor de búsqueda conecta con internet y la capacidad de programar conecta con toda la automatización del mundo digital
\end{enumerate}
\end{importantbox}

\subsection{La generalidad es la capacidad que más le falta al Agent}

\begin{importantbox}{Del Agent general a las aplicaciones verticales}
Si el Agent tuviera una capacidad de generalización suficientemente buena, no harían falta tantos Agents verticales:
\begin{itemize}
    \item Un Agent general solo necesitaría conectarse a distintas herramientas (bases de datos, API o interfaces documentales a la medida) para completar tareas de un dominio vertical
    \item Pero justamente esa capacidad de generalización es lo que más le falta al Agent actual: si puede generalizar a herramientas y entornos que no vio durante el entrenamiento
    \item El sobreajuste a los Benchmark es un riesgo: subir la puntuación en unos cuantos Benchmark no representa la capacidad de generalización en escenarios reales
\end{itemize}
\end{importantbox}

\begin{knowledgebox}{El propósito del Agent no es simular a las personas}
Yang Zhilin (杨植麟) señala una distinción importante: el propósito del Agent es la generalidad (General Purpose), no simular a las personas. Los seres humanos también son generales (Universal Constructor); que la forma de actuar del Agent se parezca a la humana es solo una coincidencia, del mismo modo que el avión se diseñó como medio de transporte, no para volar como un ave.
\end{knowledgebox}

\subsection{Resumen de la sección}

Las dos características centrales del Agent son la interacción multiturno y el uso de herramientas. El mayor reto actual es la capacidad de generalización: si puede desempeñarse bien en herramientas y entornos que no ha visto. El objetivo del Agent es la inteligencia general, no simular la conducta humana.

